%!TEX root = ../main.tex
% =============================================================
%  CAPÍTULO 8 — IMPLICANCIAS, CONCLUSIONES Y RECOMENDACIONES
%                                       (Entrega 8 · 5 % · 20/10)
% =============================================================

\chapter{Implicancias, conclusiones y recomendaciones}\label{cap:conclusiones}

\begin{guia}[Guía de la cátedra: qué se evalúa en este capítulo]
\begin{enumerate}
    \item \textbf{Respuesta explícita a la hipótesis} del
          capítulo~\ref{cap:introduccion}, con matices y contraejemplos
          si corresponde (``afirmativa, con matices'' es una respuesta
          válida; ``se cumplieron todos los objetivos'' sin evidencia no).
    \item \textbf{Implicancias}: técnicas, sociales y económicas.
    \item \textbf{Limitaciones honestas}, por categoría (técnicas, de datos,
          de validación, de mercado), sin maquillaje. Declarar como
          preliminar lo que es preliminar.
    \item \textbf{Recomendaciones} accionables y \textbf{trabajo futuro}.
    \item \textbf{Contribución al conocimiento}: qué sabe ahora la disciplina
          que antes no sabía.
\end{enumerate}
Cierra el arco narrativo: cada objetivo específico se retoma y se
indica en qué medida se alcanzó, con referencia a la sección que lo
demuestra.
\end{guia}

\section{Implicancias de los resultados}

La implicancia central del Capítulo 7 es metodológica antes que técnica:
la diferencia entre 0,368 y 0,3032 en F1 macro -\/- la medida que resume
en un solo número qué tan bien el modelo distingue los niveles de
gravedad de un siniestro, sin inflarse por acertar solo en el más
frecuente -\/- no es un detalle de medición: cambia según cómo se
separan los datos para medirla. La primera cifra sale de repartirlos al
azar entre entrenamiento y prueba; la segunda, de entrenar con años
pasados y probar contra un año posterior, como pasaría en el uso real.
Esa diferencia es la demostración de que el protocolo de evaluación
cambia la conclusión que se puede sacar de un mismo modelo. Un informe
que solo hubiera reportado la primera cifra habría sobrestimado en unos
6,5 puntos la capacidad real del sistema de generalizar a datos nuevos.
Esa brecha, medida y declarada en vez de ocultada, es en sí misma un
resultado del proyecto, no un defecto a disculpar.

Sobre esa base, una segunda medición agrega una capa más a la misma
implicancia. Evaluar el modelo contra dos años distintos, 2023→2024 y
2024→2025, dio F1 macro de 0,2859 y 0,3689 respectivamente, con
intervalos de confianza al 95\% que no se superponen. La diferencia
entre evaluar contra un año u otro es real, no ruido de muestreo. Esto
no invalida el modelo, pero sí invalida la idea de que exista un único
número de desempeño que describa al sistema de forma estable en el
tiempo: el desempeño del modelo depende de qué año se use para medirlo.
Una segunda corrida descartó una de las tres hipótesis abiertas sobre
esa dependencia: someter el entrenamiento más grande (15.559 filas de
\textless2024) al mismo tamaño que el más chico (11.644 filas de
\textless2023) no empeora el resultado sobre 2025, así que el tamaño del
conjunto de entrenamiento no explica la brecha. Esa misma corrida mostró
además que la brecha no golpea parejo entre clases: el desempeño en los
casos mortales es prácticamente estable entre ambas evaluaciones (0,126
contra 0,110), y el salto de F1 macro lo explica sobre todo un mejor
reconocimiento de los casos leves, no un cambio en las clases que más
importan para una alerta preventiva. Queda una pista nueva sin seguir
todavía: la proporción de edad sin registrar en el año evaluado pasa de
12\% en 2024 a 41\% en 2025, con 2024 marcadamente más completo que el
resto de la serie histórica, lo que sugiere que parte de la diferencia
podría deberse a una variación en la calidad del dato oficial año a año,
no solo a un cambio real en el patrón vial.

La tabla que compara, para cada nivel de gravedad, lo que el modelo
predijo contra lo que pasó en la realidad arroja una precisión de 0,084
en los siniestros graves y 0,072 en los mortales -\/- de cada 100 casos
que el modelo marca en esa clase, esa es la proporción que realmente lo
es -\/- contra una capacidad de detección de 0,600 y 0,512 -\/- de cada
100 casos reales de esa clase, esa proporción el modelo efectivamente
identifica. Eso tiene una implicancia de uso directa: el score no puede
leerse como una probabilidad individual de que a una persona le pase
algo en una zona puntual. Puede leerse, en cambio, como un criterio de
priorización relativa entre zonas, que es exactamente lo que la interfaz
hace al ordenar las zonas por nivel de riesgo en vez de exhibir el
número crudo del modelo. El trabajo de calibración de los cortes que
separan los niveles de riesgo mostrados al usuario -\/- qué proporción
de zonas, ordenadas de menor a mayor score, queda en cada nivel -\/-
sostiene esa decisión con evidencia: el corte vigente mantiene casi
plana la capacidad de detectar los casos mortales en el rango donde más
importa, a costa de detectar menos casos graves, una decisión de costo
por clase ya tomada y documentada, no una omisión.

El análisis de estabilidad termina de precisar qué tan sólido es ese
ordenamiento zona por zona: para medir cuánto cambiaría el resultado
solo por azar, se entrenaron diez modelos sobre distintas muestras
tomadas al azar (con reemplazo) del mismo conjunto de datos, con la
misma configuración de producción. El voto mayoritario de esos diez
modelos, agregado sobre 336 escenarios de remuestreo, coincide con el
nivel de riesgo vigente en el 84,9\% de las 740 zonas, pero esa
coincidencia baja a 78,5-78,9\% justo en los niveles crítico y alto, los
que más le importan a un usuario que decide una ruta. La inestabilidad
no está distribuida parejo: se concentra donde más pesa. La implicancia
de diseño directa de este hallazgo -\/- declarar la incertidumbre del
propio sistema donde existe, en vez de uniformar una seguridad que los
datos no respaldan -\/- ya está aplicada: la columna
confianza\_estabilidad se agregó a hotspots\_contexto.csv y se expone en
/hotspots sin costo adicional por request, precalculada igual que los
umbrales de riesgo. nivel\_riesgo no cambia -\/- sigue siendo el output
íntegro del modelo entrenado -\/- y confianza\_estabilidad es un dato
adicional, no un reemplazo.

El agrupamiento de siniestros en zonas de riesgo (un método que junta
puntos cercanos en el mapa para formar una sola zona, en vez de
tratarlos como hechos aislados) tiene una implicancia distinta, más
auspiciosa: cambiar el algoritmo de agrupamiento por una versión más
flexible, capaz de formar zonas de tamaño y densidad desiguales en vez
de exigir el mismo criterio en todo el mapa, no fue solo un ajuste de
parámetros: resolvió un problema real de cobertura (de 49,9\% a 75,6\%
de los siniestros de moto asignados a alguna zona) sin inflar
artificialmente el número de zonas. Que la misma metodología aplicada a
inseguridad general dé una cobertura menor (50,2\%) no es una falla del
método, es evidencia de que los delitos del GCBA están geográficamente
más dispersos que los siniestros de moto, que se repiten en las mismas
intersecciones. El patrón espacial que la hipótesis de trabajo suponía
existe, es detectable, y es distinto según la fuente de datos que se
agrupe.

La validación con usuarios reales es la implicancia que más cambia el
eje del proyecto. Antes de esa prueba, el riesgo percibido era casi
exclusivamente técnico: que el modelo no discrimine bien las clases
graves. Después de la prueba con 12 motociclistas, apareció un riesgo
distinto y comparable en magnitud: que el sistema mida algo honesto y
aun así una persona sin entrenamiento previo lo entienda mal. Dos de
doce casos siguieron interpretando el score como una predicción de
accidente pese a la aclaración textual existente, y solo 4 de 12
completaron sin ayuda la tarea de comparar rutas y explicar la
cobertura, la que exige entender qué mide el número y no solo leerlo. La
implicancia de fondo es que la precisión del modelo y la comprensión del
usuario son dos problemas distintos, y que resolver el primero sin el
segundo deja al sistema igual de inútil en la práctica.

Tomadas juntas, estas cinco líneas de evidencia (modelo, estabilidad del
score, clustering, comprensión de usuario, ingesta de prensa) sostienen
una misma conclusión transversal: MotoSafe GBA funciona mejor como
herramienta de ordenamiento y contexto declarado que como herramienta de
predicción. Esa distinción no es una limitación lateral del proyecto, es
su resultado central, y condiciona cómo deben leerse las conclusiones de
la sección siguiente.

\section{Conclusiones respecto de los objetivos y la hipótesis de trabajo}

El objetivo general del proyecto fue diseñar e implementar un MVP web
que permitiera consultar, visualizar y comparar información de riesgo
vial y contexto reciente para motociclistas del AMBA, con trazabilidad
de fuentes, cobertura declarada y límites de interpretación explícitos.
Ese objetivo se cumple: el sistema existe, está en operación, integra
las cinco fuentes previstas (siniestros, inseguridad, prensa, clima y
reportes colaborativos) y cada capa expone su procedencia, su antigüedad
y su cobertura en la interfaz, no solo en la documentación. Lo que el
Capítulo 7 agrega es que ese cumplimiento es honesto en un sentido
específico: declara también dónde el sistema no llega, en vez de
presentar una cobertura pareja que no existe.

Sobre los cinco objetivos específicos planteados para el proyecto, el
balance es el siguiente:

Integrar datos históricos de siniestralidad de motociclistas en CABA y
representarlos como zonas de referencia sobre un mapa interactivo.
Cumplido. Los 23.338 siniestros de moto georreferenciados están
cargados, agrupados en zonas con el método descrito antes y expuestos
como 740 zonas de riesgo con 75,6\% de cobertura, visibles en el mapa
con su período de origen (2019-2025).

Implementar un modelo de clasificación de gravedad condicionado a la
ocurrencia de un siniestro, documentando de forma explícita qué estima y
qué no estima. Cumplido, con una salvedad que el Capítulo 7 obliga a
nombrar: el modelo está implementado, documentado y la distinción entre
severidad condicionada y probabilidad de ocurrencia está clara en el
contrato técnico del sistema. Pero la validación con usuarios reales
muestra que esa documentación no elimina la confusión del lado del
usuario, ya que 2 de los 12 participantes siguieron confundiendo ambas
cosas, así que el objetivo está cumplido en el sistema y parcialmente
cumplido en la experiencia real de quien lo usa.

Incorporar capas de inseguridad general, noticias de prensa, clima y
reportes colaborativos sin equiparar su procedencia ni su cobertura con
la de los datos históricos de siniestros. Cumplido. Cada capa declara su
propia fecha de generación, su propio período de cobertura y su propia
fuente; el endpoint de cobertura por capas (/capas/cobertura) y la
leyenda del mapa muestran estos datos diferenciados, sin fusionarlos en
una cifra única.

Diseñar una interfaz cartográfica accesible, operable en escritorio y en
dispositivos móviles. Cumplido con evidencia de auditoría real: la
revisión de accesibilidad encontró y corrigió dos defectos concretos (el
modal de bienvenida no atrapaba el foco de teclado, y el fallo
silencioso del tráfico en vivo no avisaba al usuario), y quedó
verificada en los tamaños de pantalla de escritorio y celular probados.
Queda pendiente, declarado como tal, la prueba en un dispositivo físico
real: todo lo verificado hasta esta entrega fue emulación de navegador.

Construir una arquitectura reproducible, con API propia, base
PostgreSQL/PostGIS, proceso de ingesta de prensa programado y una vía de
despliegue en la nube documentada. Cumplido. El entorno fija versiones
exactas de cada dependencia, verificables contra su huella digital (un
valor que cambia si el archivo cambia, aunque sea en un solo carácter),
y el proceso de reentrenamiento es reproducible de punta a punta sobre
datos ya descargados: se verificó que produce exactamente el mismo
resultado, con la misma huella digital, en dos corridas distintas. El
despliegue en producción corre las migraciones de base de datos, la
carga de datos y el scraper con sus controles de concurrencia de forma
automática. Lo que falta es ejecutar ese mismo pipeline contra los
servicios externos reales en vivo, no contra una copia ya descargada.

La pregunta de investigación que guio el proyecto -\/- si una plataforma
que integre siniestralidad histórica, inseguridad, prensa y clima puede
mejorar la información disponible para un motociclista sin sobreestimar
su capacidad predictiva -\/- tiene, a esta altura del proyecto, una
respuesta afirmativa pero acotada. La plataforma sí mejora la
información disponible: antes de MotoSafe GBA, un motociclista que
quisiera cruzar estas fuentes tenía que hacerlo a mano, sin saber qué
antigüedad o cobertura tenía cada una. Lo que el Capítulo 7 obliga a
matizar es la segunda parte de la pregunta: el sistema evita
sobreestimar su capacidad predictiva en su documentación y en su
contrato técnico, pero la validación con usuarios muestra que esa
prevención todavía no alcanza por completo a la experiencia de uso real,
donde una porción de usuarios sigue leyendo el score como una predicción
individual.

La hipótesis de trabajo sostenía que la utilidad de una herramienta de
este tipo no estaría en predecir con precisión la ocurrencia de un
siniestro individual, sino en la trazabilidad de las fuentes que
integre, y que una interfaz que declare qué sabe, desde cuándo lo sabe y
con qué cobertura lo sabe resultaría más defendible y más útil que una
que convierta información heterogénea en un único indicador de riesgo.
La evidencia reunida hasta esta entrega es consistente con esta
hipótesis, aunque no la prueba de forma comparativa: los doce casos de
la validación con motociclistas evaluaron la interfaz de MotoSafe GBA
tal como se presenta, no contra una alternativa real de un único score,
así que la superioridad frente a esa alternativa sigue siendo una
hipótesis plausible, no un resultado comprobado. Lo que sí muestra la
evidencia es que el modelo, por sí solo, tiene una capacidad de
discriminación baja en las clases más graves (0,3032 en la medida
explicada antes, bajo la evaluación que respeta el orden temporal); lo
que distingue a MotoSafe GBA de una alternativa que mostrara un único
score sin contexto es precisamente que no se apoya solo en esa cifra,
sino en la combinación de procedencia, antigüedad y cobertura declaradas
capa por capa. La hipótesis no se sostiene, en cambio, de forma
automática: la validación con usuarios muestra que declarar esa
información en el sistema no garantiza por sí sola que el usuario la
incorpore a su interpretación, lo que convierte la claridad de esa
comunicación en un trabajo de diseño pendiente, no en un hecho ya
resuelto.

\section{Recomendaciones técnicas}

Las recomendaciones de esta sección parten del registro de tareas
pendientes que el proyecto ya mantiene relevado y priorizado, y
seleccionan lo que condiciona de forma directa la entrega final.

Confirmar la corrida de integración continua en GitHub, pendiente en una
entrega anterior por faltar la autorización de la cuenta de GitHub
utilizada. Ya resuelta: el flujo corrió contra la plataforma real, no
solo en local, con el job de backend (sintaxis más la suite completa de
pruebas contra PostgreSQL/PostGIS real y desechable) y el de frontend,
los dos en verde.

La pista de calidad de dato por año en la inestabilidad temporal entre
evaluaciones de distintos períodos ya se siguió hasta el final. La
anomalía de completitud que se había detectado no aparece en los campos
que describen el hecho en sí -\/- dónde ocurrió, con qué otro vehículo,
en qué tipo de vía -\/-, que se mantienen estables entre años. Aparece,
en cambio, en los dos campos que describen a la víctima, sexo y edad,
que varían juntos y de forma marcada de un año a otro. Reentrenar sin el
sexo de la víctima como variable del modelo confirma que esa variable
explica la mayor parte de la inestabilidad: la brecha entre evaluar
contra distintos años prácticamente desaparece, y además mejora el
reconocimiento de los casos leves. Pero no es una mejora sin costo: en
el período más reciente evaluado, empeora el reconocimiento de los casos
graves y mortales, que son los que más importan para una alerta
preventiva. Con ese costo cuantificado, la decisión de sacar o no esa
variable del modelo queda planteada para resolverse antes de la entrega
final, no tomada todavía.

Seguir separando incrementalmente el archivo principal de la interfaz
del frontend. El componente concentraba 2140 líneas de código y
alrededor de 30 bloques que recalculan datos cuando algo cambia en
pantalla; una medición real ya descartó que esto haga que la pantalla se
vuelva a dibujar de más, así que la urgencia no es de velocidad sino de
mantenimiento: un archivo de ese tamaño es más difícil de revisar y de
modificar sin introducir un error de alcance no deseado. Se extrajeron
nueve módulos autocontenidos con una interfaz de props chica (el panel
de capas, el panel de opciones de ruta, el menú de herramientas, la hoja
de reporte, la tarjeta de maniobra activa, el resumen de navegación con
velocímetro, el panel de detalle de zona, la leyenda de riesgo e
inseguridad, y el buscador de lugar/calle), sin cambio de
comportamiento, bajando el archivo a 1959 líneas en su punto más bajo
registrado. Trabajo posterior (el sistema de usuarios y otras
correcciones) volvió a sumar código: al cierre de esta entrega, App.jsx
tiene 2099 líneas. La revisión mostró que esa dependencia cruzada vivía
sobre todo en los cálculos (los bloques que calculan velocidad, tiempo
estimado de llegada y paso activo durante la navegación), no en el
código que solo define qué se muestra en pantalla: una vez separado el
cálculo de la presentación visual, los bloques resultaron más fáciles de
separar de lo previsto. Lo que queda -\/- el formulario de ruta con
origen y destino, el panel de actividad reciente y los banners de error
de capas -\/- sí tiene acoplamiento cruzado real entre los campos del
formulario y necesita una sesión dedicada a desenredarlo antes de
separarlo.

Ejecutar el pipeline de datos de punta a punta contra los servicios
reales en producción. La reproducibilidad del pipeline ya está
demostrada contra una copia de datos ya descargada, con hash idéntico al
artefacto vigente. Lo que falta es correr los mismos pasos contra los
servicios en vivo (descarga de delitos, clima y corredores desde sus
fuentes originales), no contra una copia local, para confirmar que el
proceso de actualización periódica funciona en las condiciones reales de
operación y no solo en un entorno de prueba controlado.

Construir un conjunto etiquetado para la extracción de calles del
scraper de prensa. La deduplicación semántica de notas publicadas por
medios distintos sobre el mismo hecho se mide de forma automática por
ventana de días, radio geográfico y similitud de título, sin revisión
humana y sin descartar ninguna fila. Medida contra 450 filas reales
re-geocodificadas (426 a nivel de barrio, 24 a nivel de calle o
intersección): los nueve duplicados reales de la muestra se detectan
igual, porque el umbral de similitud de título por sí solo ya alcanza
cuando dos notas comparten barrio, sin necesitar coordenada precisa. El
hueco real medido es más angosto: notas de barrio con títulos
editorialmente distintos del mismo hecho, caso que no aparece ninguna
vez en esas 450 filas y que se detectó una sola vez en una corrida en
vivo (Villa Pueyrredón, cuatro medios). Se acepta ese residuo como
riesgo conocido y acotado en vez de perseguirlo con una heurística sin
sustento. Lo que sigue pendiente es otra cosa: la extracción de calles e
intersecciones del texto de las notas nunca se validó contra un conjunto
etiquetado a mano, solo contra la inspección manual de casos puntuales.

Decidir el destino de las filas de prensa de la fuente retirada. Las
notas provenientes de Google Noticias no se borraron cuando esa fuente
dejó de usarse como origen de datos: siguen en la base, pero la API las
excluye de inmediato (no se muestran ni se cuentan), igual que las notas
marcadas como duplicado semántico de otra. Es una decisión de criterio
pendiente, no una tarea técnica compleja: definir si esas filas deben
purgarse antes de la entrega final o si conservarlas en la base sin
mostrarlas de forma indefinida es un criterio aceptable.

\section{Recomendaciones de adopción, gobernanza y uso}

Simplificar la pantalla principal por defecto, tal como lo pidió la
validación con usuarios reales. Ya aplicada: la cobertura por fuente y
el desglose de puntos de cada tarjeta de ruta quedaron detrás de un
control opcional ("Ver cómo se calcula"), y por default solo se muestra
tiempo, distancia y la advertencia principal, que es exactamente lo que
pidieron los usuarios de la prueba real. Falta el paso que cierra el
ciclo: repetir la validación con usuarios reales para confirmar que la
fricción medida anteriormente efectivamente bajó.

Rediseñar la explicación visual del score y de la cobertura. Aplicada
parcialmente: el texto aclaratorio dependía de un tooltip de mouse,
inútil en un celular usado en la calle, que es el contexto real de uso.
Se reemplazó por un ejemplo visual siempre disponible con números reales
("cobertura de siniestros 40\% significa 4 de cada 10 tramos..."),
visible junto al mismo control de detalle agregado en la recomendación
anterior. Queda parcial, no cerrada, porque sigue sin confirmarse con
una nueva ronda de usuarios si ese ejemplo efectivamente reduce la
confusión que hoy afecta a uno de cada seis casos.

Corregir la nomenclatura de los filtros de seguridad. Ya aplicada,
corrigiendo los dos defectos puntuales que señaló la prueba con
usuarios: el filtro de prensa llamado "Seguridad" se renombró a
"Violencia", que es lo que esa categoría realmente trae según las
palabras clave del scraper (tiroteos, homicidios, secuestros, ya
separados de robo y de choque); y la nota que explica que apagar la capa
de inseguridad no la saca del cálculo de ruta, que ya existía pero
genérica al pie del panel, se duplicó en una versión más corta pegada
directamente a esas opciones.

Formalizar una política de privacidad para los datos de ubicación de los
reportes colaborativos. Ya resuelta con un documento nuevo en la carpeta
de operación del repositorio, que describe qué guarda cada reporte
(latitud y longitud exactas, categoría, un identificador de dispositivo
sin dato personal), qué no guarda (IP persistida, nombre, email), cómo
se retiene y se borra, y declara de forma explícita una limitación que
antes no estaba escrita en ningún lado: la ubicación se guarda sin
redondear, con el costo de privacidad real que eso implica si alguien
correlacionara varios reportes del mismo dispositivo.

Validar el diferencial de mercado antes de pensar en escalar. MotoSafe
GBA no compite en un vacío: cualquier aplicación de navegación ya
instalada podría sumar una función similar. El diferencial real del
proyecto, el foco hiperlocal en motociclistas del GBA con reportes
comunitarios, corredores de riesgo y prensa georreferenciada combinados
en un mismo mapa, todavía no se validó contra una base de usuarios que
lo sostenga más allá de la prueba de doce casos. Antes de evaluar
cualquier escalamiento posterior a la entrega académica, conviene medir
si ese diferencial efectivamente retiene usuarios con el tiempo, no solo
si les resulta útil en una sesión de prueba moderada.

Sostener la gobernanza de un solo responsable técnico con registro
explícito. El desarrollo a cargo de una sola persona no es, en sí mismo,
un riesgo no gestionado: el registro de decisiones y de cada paso en la
bitácora de desarrollo cumple la función que en un equipo más grande
cumpliría un proceso formal de revisión. La recomendación, de cara a una
eventual continuidad del proyecto más allá de la tesis, es que esa misma
disciplina de registro se mantenga si se suma una segunda persona al
desarrollo, en vez de diluirse con el crecimiento del equipo.

\section{Cierre}

Un sistema que predijera con precisión no necesitaría declarar su
incertidumbre por zona, ni su cobertura por capa, ni la diferencia entre
dos protocolos de evaluación: alcanzaría con mostrar un número. MotoSafe
GBA hace lo contrario porque no tiene otra opción honesta: el modelo
discrimina poco en las clases más graves, el acuerdo entre las muestras
tomadas al azar para medir su estabilidad cae justo donde más importa, y
la propia validación con usuarios reales muestra que ni siquiera una
aclaración textual correcta garantiza que se entienda bien. Declarar
estos límites, en vez de ocultarlos detrás de una interfaz prolija, es
la decisión de diseño que atraviesa todo el proyecto.

Las conclusiones anteriores muestran que el objetivo general y los cinco
objetivos específicos se cumplen, con dos salvedades declaradas y no
escondidas: la comprensión del usuario todavía no está a la altura de la
honestidad técnica del sistema, y la ejecución repetida y sostenida del
pipeline de reentrenamiento contra los servicios externos en vivo
(delitos, clima, geometría de corredores) sigue sin confirmarse más allá
de la corrida puntual ya hecha (detalle en 8.3); el scraper de prensa sí
corre así, de forma repetida y verificada, desde GitHub Actions. Las
recomendaciones de las dos secciones anteriores no son una lista de
deseos para un proyecto futuro: son los pasos concretos y ya
identificados que separan a MotoSafe GBA, en su estado actual, de una
versión que cierre esas dos brechas antes de la entrega final. El
criterio que sostuvo el proyecto hasta esta entrega, medir antes de
afirmar y declarar el límite en vez de disimularlo, es el mismo que
debería guiar el trabajo que queda.
